\documentclass[reqno]{amsart}
\usepackage[T1]{fontenc}
\usepackage{amsmath,amssymb,amsthm}
\usepackage{xr-hyper}
\usepackage[hidelinks]{hyperref}
\externaldocument[da-]{20261126T030300--the-data__workbook}
\externaldocument[ob-]{20261126T040100--spaces-forms-and-twists__workbook}
\externaldocument[cg-]{20261125T014100--congruence-modules__workbook}
\externaldocument[pc-]{20261125T014700--pseudocharacters-and-the-length-bound__workbook}
\externaldocument[gr-]{20261125T011000--galois-representations-black-box__workbook}
\externaldocument[tt-]{20261125T004100--tate-twists-and-sha__workbook}
\externaldocument[dp-]{20261126T030900--denominators-of-projections-to-lifts__workbook}
\externaldocument[n3-]{20261126T051800--the-local-loss-and-the-sparse-places__workbook}
\externaldocument[cp-]{20261126T070000--clearing-a-projection-two-classes__workbook}

\newcounter{chap}
% Chapter numbers, generated from 20261125T012000--grh-workbook__guide.org
% by scripts/make-chapters.py.  Do not edit: change the order of the
% index and rerun the script (./preview-site.sh and CI do this).
%
% Usage in a sheet, after \newcounter{chap}:
%   % Chapter numbers, generated from 20261125T012000--grh-workbook__guide.org
% by scripts/make-chapters.py.  Do not edit: change the order of the
% index and rerun the script (./preview-site.sh and CI do this).
%
% Usage in a sheet, after \newcounter{chap}:
%   % Chapter numbers, generated from 20261125T012000--grh-workbook__guide.org
% by scripts/make-chapters.py.  Do not edit: change the order of the
% index and rerun the script (./preview-site.sh and CI do this).
%
% Usage in a sheet, after \newcounter{chap}:
%   \input{chapters}\chapnumber{<sheet ID>}
\makeatletter
\providecommand\chapters@set[2]{\@namedef{chapters@#1}{#2}}
\chapters@set{20261125T001000}{1}% zeros-to-ramp-sums
\chapters@set{20261125T001100}{2}% inert-primes-and-estimate-b
\chapters@set{20261125T005000}{3}% galois-fourier-and-averaging
\chapters@set{20261125T005500}{4}% columns-from-a-sphere
\chapters@set{20261125T000000}{5}% slope-method
\chapters@set{20261126T030100}{6}% a-twist-for-each-prime
\chapters@set{20261125T003000}{7}% the-chain-and-the-square-root
\chapters@set{20261125T010000}{8}% rank-lemma-and-zero-estimate
\chapters@set{20261125T010500}{9}% poincare-series-and-near-orthogonality
\chapters@set{20261124T223000}{10}% elliptic-curves-for-the-proof
\chapters@set{20261125T009500}{11}% fock-model-and-theta-intertwiners
\chapters@set{20261124T223100}{12}% complex-multiplication-and-inert-reduction
\chapters@set{20261125T008000}{13}% one-cm-determinant
\chapters@set{20261124T214500}{14}% supersingular-theta-divisibility
\chapters@set{20261124T214600}{15}% the-isotropic-plane
\chapters@set{20261125T013700}{16}% pairings-on-a-formal-group-of-height-two
\chapters@set{20261125T013500}{17}% hecke-nondegeneracy-on-the-tree
\chapters@set{20261126T030200}{18}% the-bound-is-attained
\chapters@set{20261125T009000}{19}% kronecker-limit-and-chowla-selberg
\chapters@set{20261125T002000}{20}% definite-theta-lifts
\chapters@set{20261125T015500}{21}% siegel-weil-in-the-smallest-case
\chapters@set{20261125T002100}{22}% definite-theta-lifts-and-norms
\chapters@set{20261125T016500}{23}% hecke-operators-and-theta-lifts
\chapters@set{20261125T016700}{24}% divisibility-along-the-chain
\chapters@set{20261125T004000}{25}% denominators-and-sha
\chapters@set{20261125T004100}{26}% tate-twists-and-sha
\chapters@set{20261126T020430}{27}% the-minus-part-by-gauss-sums
\chapters@set{20261125T014100}{28}% congruence-modules
\chapters@set{20261125T014500}{29}% eisenstein-congruences-and-selmer-classes
\chapters@set{20261125T014700}{30}% pseudocharacters-and-the-length-bound
\chapters@set{20261125T011000}{31}% galois-representations-black-box
\chapters@set{20261126T030300}{32}% the-data
\chapters@set{20261126T020210}{33}% bessel-poincare-series-and-their-gram-operators
\chapters@set{20261126T030400}{34}% the-seesaw-identity
\chapters@set{20261126T030500}{35}% divisibility-of-the-composite-lift
\chapters@set{20261126T030600}{36}% theorem-a-short-forms
\chapters@set{20261126T030700}{37}% the-determinant-of-a-chain
\chapters@set{20261126T030800}{38}% the-cost-of-one-lift
\chapters@set{20261126T030900}{39}% denominators-of-projections-to-lifts
\chapters@set{20261126T031000}{40}% twists-with-few-congruences
\chapters@set{20261126T031100}{41}% theorem-c-and-the-main-inequality
\providecommand\chapnumber[1]{%
  \@ifundefined{chapters@#1}%
    {\PackageError{chapters}{Sheet #1 is not listed in the index}%
      {List it in 20261125T012000--grh-workbook__guide.org\MessageBreak
       and run scripts/make-chapters.py.}}%
    {\edef\thechap{\@nameuse{chapters@#1}}}}
\makeatother
\chapnumber{<sheet ID>}
\makeatletter
\providecommand\chapters@set[2]{\@namedef{chapters@#1}{#2}}
\chapters@set{20261125T001000}{1}% zeros-to-ramp-sums
\chapters@set{20261125T001100}{2}% inert-primes-and-estimate-b
\chapters@set{20261125T005000}{3}% galois-fourier-and-averaging
\chapters@set{20261125T005500}{4}% columns-from-a-sphere
\chapters@set{20261125T000000}{5}% slope-method
\chapters@set{20261126T030100}{6}% a-twist-for-each-prime
\chapters@set{20261125T003000}{7}% the-chain-and-the-square-root
\chapters@set{20261125T010000}{8}% rank-lemma-and-zero-estimate
\chapters@set{20261125T010500}{9}% poincare-series-and-near-orthogonality
\chapters@set{20261124T223000}{10}% elliptic-curves-for-the-proof
\chapters@set{20261125T009500}{11}% fock-model-and-theta-intertwiners
\chapters@set{20261124T223100}{12}% complex-multiplication-and-inert-reduction
\chapters@set{20261125T008000}{13}% one-cm-determinant
\chapters@set{20261124T214500}{14}% supersingular-theta-divisibility
\chapters@set{20261124T214600}{15}% the-isotropic-plane
\chapters@set{20261125T013700}{16}% pairings-on-a-formal-group-of-height-two
\chapters@set{20261125T013500}{17}% hecke-nondegeneracy-on-the-tree
\chapters@set{20261126T030200}{18}% the-bound-is-attained
\chapters@set{20261125T009000}{19}% kronecker-limit-and-chowla-selberg
\chapters@set{20261125T002000}{20}% definite-theta-lifts
\chapters@set{20261125T015500}{21}% siegel-weil-in-the-smallest-case
\chapters@set{20261125T002100}{22}% definite-theta-lifts-and-norms
\chapters@set{20261125T016500}{23}% hecke-operators-and-theta-lifts
\chapters@set{20261125T016700}{24}% divisibility-along-the-chain
\chapters@set{20261125T004000}{25}% denominators-and-sha
\chapters@set{20261125T004100}{26}% tate-twists-and-sha
\chapters@set{20261126T020430}{27}% the-minus-part-by-gauss-sums
\chapters@set{20261125T014100}{28}% congruence-modules
\chapters@set{20261125T014500}{29}% eisenstein-congruences-and-selmer-classes
\chapters@set{20261125T014700}{30}% pseudocharacters-and-the-length-bound
\chapters@set{20261125T011000}{31}% galois-representations-black-box
\chapters@set{20261126T030300}{32}% the-data
\chapters@set{20261126T020210}{33}% bessel-poincare-series-and-their-gram-operators
\chapters@set{20261126T030400}{34}% the-seesaw-identity
\chapters@set{20261126T030500}{35}% divisibility-of-the-composite-lift
\chapters@set{20261126T030600}{36}% theorem-a-short-forms
\chapters@set{20261126T030700}{37}% the-determinant-of-a-chain
\chapters@set{20261126T030800}{38}% the-cost-of-one-lift
\chapters@set{20261126T030900}{39}% denominators-of-projections-to-lifts
\chapters@set{20261126T031000}{40}% twists-with-few-congruences
\chapters@set{20261126T031100}{41}% theorem-c-and-the-main-inequality
\providecommand\chapnumber[1]{%
  \@ifundefined{chapters@#1}%
    {\PackageError{chapters}{Sheet #1 is not listed in the index}%
      {List it in 20261125T012000--grh-workbook__guide.org\MessageBreak
       and run scripts/make-chapters.py.}}%
    {\edef\thechap{\@nameuse{chapters@#1}}}}
\makeatother
\chapnumber{<sheet ID>}
\makeatletter
\providecommand\chapters@set[2]{\@namedef{chapters@#1}{#2}}
\chapters@set{20261125T001000}{1}% zeros-to-ramp-sums
\chapters@set{20261125T001100}{2}% inert-primes-and-estimate-b
\chapters@set{20261125T005000}{3}% galois-fourier-and-averaging
\chapters@set{20261125T005500}{4}% columns-from-a-sphere
\chapters@set{20261125T000000}{5}% slope-method
\chapters@set{20261126T030100}{6}% a-twist-for-each-prime
\chapters@set{20261125T003000}{7}% the-chain-and-the-square-root
\chapters@set{20261125T010000}{8}% rank-lemma-and-zero-estimate
\chapters@set{20261125T010500}{9}% poincare-series-and-near-orthogonality
\chapters@set{20261124T223000}{10}% elliptic-curves-for-the-proof
\chapters@set{20261125T009500}{11}% fock-model-and-theta-intertwiners
\chapters@set{20261124T223100}{12}% complex-multiplication-and-inert-reduction
\chapters@set{20261125T008000}{13}% one-cm-determinant
\chapters@set{20261124T214500}{14}% supersingular-theta-divisibility
\chapters@set{20261124T214600}{15}% the-isotropic-plane
\chapters@set{20261125T013700}{16}% pairings-on-a-formal-group-of-height-two
\chapters@set{20261125T013500}{17}% hecke-nondegeneracy-on-the-tree
\chapters@set{20261126T030200}{18}% the-bound-is-attained
\chapters@set{20261125T009000}{19}% kronecker-limit-and-chowla-selberg
\chapters@set{20261125T002000}{20}% definite-theta-lifts
\chapters@set{20261125T015500}{21}% siegel-weil-in-the-smallest-case
\chapters@set{20261125T002100}{22}% definite-theta-lifts-and-norms
\chapters@set{20261125T016500}{23}% hecke-operators-and-theta-lifts
\chapters@set{20261125T016700}{24}% divisibility-along-the-chain
\chapters@set{20261125T004000}{25}% denominators-and-sha
\chapters@set{20261125T004100}{26}% tate-twists-and-sha
\chapters@set{20261126T020430}{27}% the-minus-part-by-gauss-sums
\chapters@set{20261125T014100}{28}% congruence-modules
\chapters@set{20261125T014500}{29}% eisenstein-congruences-and-selmer-classes
\chapters@set{20261125T014700}{30}% pseudocharacters-and-the-length-bound
\chapters@set{20261125T011000}{31}% galois-representations-black-box
\chapters@set{20261126T030300}{32}% the-data
\chapters@set{20261126T020210}{33}% bessel-poincare-series-and-their-gram-operators
\chapters@set{20261126T030400}{34}% the-seesaw-identity
\chapters@set{20261126T030500}{35}% divisibility-of-the-composite-lift
\chapters@set{20261126T030600}{36}% theorem-a-short-forms
\chapters@set{20261126T030700}{37}% the-determinant-of-a-chain
\chapters@set{20261126T030800}{38}% the-cost-of-one-lift
\chapters@set{20261126T030900}{39}% denominators-of-projections-to-lifts
\chapters@set{20261126T031000}{40}% twists-with-few-congruences
\chapters@set{20261126T031100}{41}% theorem-c-and-the-main-inequality
\providecommand\chapnumber[1]{%
  \@ifundefined{chapters@#1}%
    {\PackageError{chapters}{Sheet #1 is not listed in the index}%
      {List it in 20261125T012000--grh-workbook__guide.org\MessageBreak
       and run scripts/make-chapters.py.}}%
    {\edef\thechap{\@nameuse{chapters@#1}}}}
\makeatother

\ifcsname chapters@20261126T060400\endcsname
  \chapnumber{20261126T060400}
\else
  \renewcommand{\thechap}{?}
\fi
\renewcommand{\thesection}{\thechap.\arabic{section}}
\theoremstyle{definition}
\newtheorem{problem}{Problem}[section]
\newtheorem{definition}{Definition}[section]
\theoremstyle{plain}
\newtheorem{theorem}{Theorem}[section]
\newcommand{\Z}{\mathbb{Z}}
\newcommand{\Q}{\mathbb{Q}}
\newcommand{\C}{\mathbb{C}}
\newcommand{\F}{\mathbb{F}}
\newcommand{\Oc}{\mathcal{O}}
\newcommand{\val}{\operatorname{val}}
\newcommand{\ord}{\operatorname{ord}}
\newcommand{\len}{\operatorname{len}}
\newcommand{\tr}{\operatorname{tr}}
\newcommand{\Fitt}{\mathrm{Fitt}_0}
\newcommand{\GL}{\mathrm{GL}}
\newcommand{\cyc}{\chi_{\mathrm{cyc}}}
\newcommand{\DP}{\mathfrak{O}}
\newcommand{\cl}{\mathrm{cl}}
\newcommand{\loc}{\mathrm{loc}}
\newcommand{\sep}{c_{\mathrm{sep}}}
\newcommand{\Lc}{\mathcal{L}}
\newcommand{\Tb}{\mathbb{T}}
\newcommand{\half}{(\Delta/2)}
\setlength{\emergencystretch}{3em}

\title{Congruences bounded by lengths}
\author{}
\date{November 26, 2026}

\begin{document}
\maketitle
\refstepcounter{chap}\label{chap}

\paragraph{Prerequisites.}
Discrete valuation rings, lengths, elementary divisors; Fitting ideals \cite[\S20.2]{Ei}.  Representations of \(S_3\).  From Chapter~\ref{pc-chap}: Section~\ref{pc-sec:relations}, Problem~\ref{pc-pr:graph-map} and Theorem~\ref{pc-thm:graph}; Problem~\ref{pc-pr:sha-lower}.  Problem~\ref{cg-pr:toy}.  From Chapter~\ref{dp-chap}: Section~\ref{dp-sec:statement} with all its definitions, Theorem~\ref{dp-thm:denominators}, Problems~\ref{dp-pr:budget} and~\ref{dp-pr:det-bound}.  From Chapter~\ref{cp-chap}: Problems~\ref{cp-pr:generator}, \ref{cp-pr:gl-order} and~\ref{cp-pr:pairing}.  The local loss \(\loc_\ell\) of Definition~\ref{n3-def:loss}; Theorem~\ref{tt-thm:oddsha} for \(X_{\mathrm{odd}}\).

\medskip
Theorem~\ref{dp-thm:denominators} bounds the denominator of the projection onto a space of lifts at \(\iota\) by \(\cl_\ell\log\ell\).  This sheet serves that step: it states the bound for the eigenvalues of one clearing operator, Theorem~\ref{thm:eigenvalues}, and deduces the bound for the projection from it (Problem~\ref{pr:deduce}).  In Problem~\ref{pc-pr:sha-lower} a congruence of exponent \(e\) gave a cyclic module of classes of length \(e\).  Here bounds for the lengths of modules of cocycles enter, through a Fitting ideal, an upper bound for the order at the lift of a coefficient of a Hecke element.  The element multiplies certain determinant families, which vanish on the lifts, into a module of traces; in the two-system model it can be chosen to vanish at the other system.  How small can its order at the lift be?  We work this out on \(S_3\) at \(3\), where the congruence exponent and the length of the normalized cocycles both equal \(1\), and then count the costs of separating characters and of tame ramification in the setting of Chapter~\ref{dp-chap}.

\section{The smallest example}\label{sec:model}

Let \(G = S_3\) and \(\Oc = \Z_3\).  Let \(U_1(g) := \mathrm{diag}(\mathrm{sgn}(g), 1)\), and let \(U_2(g)\) be the matrix of \(g\) on the lattice \(\{(a,b,c) \in \Z_3^3 : a+b+c = 0\}\), permuting coordinates, in the basis \(e_1 - e_2\), \(e_2 - e_3\).  Think of \(U_1\) as a lift (a block \(\mathrm{sgn}\) plus the trivial character) and of \(U_2\) as another system.  Let \(\Tb \subseteq \Z_3 \times \Z_3\) be the \(\Z_3\)-algebra generated by the pairs of coefficients of the characteristic polynomials of \(U_1(g)\) and \(U_2(g)\), \(g \in G\): a toy Hecke algebra (in this sheet \(\Tb\) is this algebra, not \(\mathbb T^{\mathcal M}\)).

\begin{problem}\label{pr:model-traces}
The pairs \((\tr U_1(g), \tr U_2(g))\) are \((2,2)\), \((2,-1)\), \((0,0)\) for \(g = 1\), a \(3\)-cycle, a transposition.  The characteristic polynomials of \(U_1(g)\) and \(U_2(g)\) agree modulo \(3\) for every \(g\), and \(\Tb = \{(a,b) : a \equiv b \pmod 3\}\).
\end{problem}

\begin{problem}\label{pr:model-congruence}
The elements of \(\Tb\) that vanish at \(U_2\) are the \((3u, 0)\), \(u \in \Z_3\); so the congruence between \(U_1\) and \(U_2\) has exponent \(1\) (compare Problem~\ref{cg-pr:toy} with \(N = 3\)).  The representation \(U_2 \otimes \Q_3\) is irreducible, and the reduction of \(U_2\) modulo \(3\) has semisimplification \(1 \oplus \mathrm{sgn}\).
\end{problem}

\section{Ribet cocycles}\label{sec:ribet}

Let \(\Oc\) be a commutative ring, \(G\) a group, and \(U_1, \dots, U_n \colon G \to \GL_q(\Oc)\) homomorphisms; write \(U(g) - 1 := (U_i(g) - 1)_i \in M_q(\Oc)^n\), and let \(\Lc\) be the \(\Oc\)-span of the \(U(g) - 1\).  In this sheet the sans-serif letters \(\mathsf B, \mathsf C, \mathsf K, \mathsf M, \mathsf N, \mathsf R\) denote modules, rings and scalars of commutative algebra, distinct from the reserved \(B, K, M, N\).

\begin{definition}\label{def:ribet}
Let \(\rho \colon G \to \GL_s(\Oc)\) be a homomorphism (a \emph{block}) and \(\mathsf B\) an \(\Oc\)-module.  A \emph{Ribet cocycle} for \(\rho\) with values in \(\mathsf B^s\) is a map \(z \colon G \to \mathsf B^s\) with \(z(gh) = z(g) + \rho(g)z(h)\) for all \(g, h\), that obeys the relations among the \(U(g) - 1\): for every finitely supported \(\theta \colon G \to \Oc\) with \(\sum_g \theta(g)(U(g) - 1) = 0\), also \(\sum_g \theta(g) z(g) = 0\).  It is \emph{normalized at} \(g_* \in G\) if \(z(g_*) = 0\).
\end{definition}

For \(\mathsf B = \Lambda/\alpha\) these are the cocycles of Section~\ref{pc-sec:relations}.

\begin{problem}\label{pr:normalize}
Suppose \(\rho(g_*) = y \cdot 1\) with \(y - 1 \in \Oc^\times\).  Then a coboundary \(g \mapsto \rho(g)m - m\) that vanishes at \(g_*\) is \(0\), and every Ribet cocycle for \(\rho\) differs by exactly one coboundary from a Ribet cocycle normalized at \(g_*\).
\end{problem}

So normalized cocycles stand for classes.  Return to the model of Section~\ref{sec:model}, with \(U = (U_1, U_2)\), the block \(\rho = \mathrm{sgn}\), and \(g_* = \tau := (1\,2)\), where \(\mathrm{sgn}(\tau) - 1 = -2\) is a unit.

\begin{problem}\label{pr:model-lattice}
By machine: \(\Lc \subseteq M_2(\Z_3)^2\) is free of rank \(5\).  Let \(\Lc'\) be the span of the elements \((U(g) - \mathrm{sgn}(g))(U(h) - 1)\), \(g, h \in G\), and of \(U(\tau) - 1\).  Then \(\Lc/\Lc' \cong \Z/3\); without \(U(\tau) - 1\) the quotient is \(\Z_3 \oplus \Z/3\).
\end{problem}

\begin{problem}\label{pr:model-cocycles}
For every \(n \ge 1\), the Ribet cocycles for \(\mathrm{sgn}\) with values in \(\Z/3^n\) normalized at \(\tau\) form a group of order \(3\), and all Ribet cocycles for \(\mathrm{sgn}\) a group of order \(3^{n+1}\).  With \(U = U_1\) alone, the normalized ones are \(0\).
\end{problem}

With values in \(\Q_3/\Z_3\) the normalized cocycles therefore have length \(1\), the exponent of the congruence of Problem~\ref{pr:model-congruence}; with \(U_1\) alone the module is zero.

\section{Lengths bound Fitting ideals}\label{sec:fitting}

From now on \(\Oc\) is a discrete valuation ring with uniformizer \(\varpi\), fraction field \(\mathrm{Frac}(\Oc)\) and normalized valuation \(\ord\).  For a commutative ring \(\mathsf R\) and an \(\mathsf R\)-submodule \(\mathsf N \subseteq \mathsf R^s\), \(\Fitt(\mathsf R^s/\mathsf N)\) is the ideal generated by the determinants of the \(s \times s\) matrices whose rows lie in \(\mathsf N\).

\begin{problem}\label{pr:fitt-length}
Let \(\mathsf M = \Oc^n/\mathsf N\) have finite length.  Then \(\Fitt(\mathsf M) = \varpi^{\len \mathsf M}\Oc\) and \(\len \operatorname{Hom}_{\Oc}(\mathsf M, \mathrm{Frac}(\Oc)/\Oc) = \len \mathsf M\).  For the rows \((3, 0)\), \((1, 9)\) over \(\Z_3\), \(\mathsf M \cong \Z/27\).
\end{problem}

In the model, the quotient \(\Lc/\Lc'\) of Problem~\ref{pr:model-lattice} has \(\Fitt = 3\Z_3\).  For the following problems let \(\mathsf R\) be a commutative \(\Oc\)-algebra.  Let \(\Phi \colon \Oc^n \to \mathsf R^s\) be \(\Oc\)-linear, \(p \ne 0\) in \(\Oc\), and \(\mathsf N \subseteq \mathsf R^s\) an \(\mathsf R\)-submodule with \(\Phi(p\Oc^n) \subseteq \mathsf N\).  The \emph{forms vanishing on \(\Phi^{-1}\mathsf N\)} are the \(\Oc\)-linear \(\Psi \colon \Oc^n \to \Oc/p\) with \(\Psi(v) = 0\) whenever \(\Phi(v) \in \mathsf N\).

\begin{problem}\label{pr:forms}
\(\Phi(u) \in \mathsf N\) if and only if every form vanishing on \(\Phi^{-1}\mathsf N\) kills \(u\).
\end{problem}

\begin{problem}\label{pr:fitting-step}
Suppose \(\Phi(v) \notin \mathsf N\).
\begin{enumerate}
\item[(a)] Some \(\mathbf w = \Phi(\varpi^k v)\) has \(\mathbf w \notin \mathsf N\) and \(\varpi \mathbf w \in \mathsf N\).
\item[(b)] For such \(\mathbf w\), \(\varpi\,\Fitt(\mathsf R^s/(\mathsf N + \mathsf R\mathbf w)) \subseteq \Fitt(\mathsf R^s/\mathsf N)\), and the module of forms vanishing on \(\Phi^{-1}(\mathsf N + \mathsf R\mathbf w)\) is strictly smaller than that for \(\mathsf N\).
\item[(c)] If \(\mathsf R\) is noetherian, there is \(c \in \Oc\) with \(\ord c\) at most the length of the module of forms vanishing on \(\Phi^{-1}\mathsf N\) and \(c\,\Fitt\bigl(\mathsf R^s/(\mathsf N + \mathsf R \cdot \Phi(\Oc^n))\bigr) \subseteq \Fitt(\mathsf R^s/\mathsf N)\).
\end{enumerate}
\end{problem}

The proof uses a relative form of~(c), for noetherian \(\mathsf R\).  Let \(z \colon G \to \Oc^n\) be a family of vectors, \(V\) an \(\Oc\)-module of functions \(G \to \Oc/p\), and \(F\) the module of the functions \(\Psi \circ z\) for the forms \(\Psi\) vanishing on \(\Phi^{-1}\mathsf N\).  Then there are \(c \in \Oc\) and an \(\mathsf R\)-submodule \(\mathsf K \supseteq \mathsf N\) with \(\ord c \le \len(F/(F \cap V))\), \(c\,\Fitt(\mathsf R^s/\mathsf K) \subseteq \Fitt(\mathsf R^s/\mathsf N)\), and \(\Phi(v) \in \mathsf K\) for every \(v\) killed by every \(\Oc\)-linear \(\Psi \colon \Oc^n \to \Oc/p\) such that \(\Psi \circ z \in V\).  Here is the block bound of the proof with no auxiliary indices.  Choose \(g_1', \dots, g_m' \in G\) and coordinates \(z(g) \in \Oc^m\) with \(U(g) - 1 = \sum_i z(g)_i(U(g_i') - 1)\), and write \([g]\) for \(g\) in \(\Oc[G]\).  Let \(\mathsf N_0 \subseteq \Oc[G]^m\) be a left submodule containing the rows \(z(gh) - z(g) - [g]z(h)\), \(z(g_i') - e_i\), and \(\sum_g \theta(g)z(g)\) for every finitely supported relation \(\sum_g\theta(g)(U(g) - 1) = 0\), and a row whose image under \([g] \mapsto \rho(g)\) agrees with that of \(z(g_*)\).  If \(\rho(G)\) spans \(M_s(\Oc)\) and the Ribet cocycles for \(\rho\) with values in \((\mathrm{Frac}(\Oc)/\Oc)^s\) normalized at \(g_*\) have length at most \(b_{\mathrm{main}} \in \mathbb{N}\), some \(m \times m\) matrix with rows in \(\mathsf N_0\) has, after applying \(\rho\) entrywise and flattening, determinant of order at most \(b_{\mathrm{main}}\).  With auxiliary indices the proof takes tied rows over the polynomial group algebra, uses a scalar normalizer as in Problem~\ref{pr:normalize}, imposes vanishing also at the first element of each tie (the elements \(s_{jb}\) of the general lemma in Section~\ref{sec:costs}), and sets all variables to zero before taking this determinant.  An element of a Fitting ideal acts like a clearing scalar:

\begin{problem}\label{pr:det-kills}
For a square matrix \(\mathsf B\) over a commutative ring, \(\det \mathsf B\) kills the cokernel of \(\mathsf B\).
\end{problem}

\begin{definition}\label{def:weak}
A homomorphism \(U \colon G \to \GL_q(\Oc)\) is a \emph{weak lift} if \(\det(1 - U(g)) = 0\) for every \(g \in G\).
\end{definition}

For a system of forms the proof asks this of the characteristic polynomials of its normalized representation at every \(g \in G_E\).  For an attached system satisfying the lift trace formula of Definition~\ref{dp-def:main-summand}, the normalized representation is a weak lift, since its trace at every group element includes the trivial character \(\cyc^0\) among the stated summands.

\begin{problem}\label{pr:weak}
If \(U = \rho \oplus 1\) in some basis, then \(U\) is a weak lift and \(\det \sum_g \omega(g)(U(g) - 1) = 0\) for every finitely supported \(\omega \colon G \to \Oc\).  Conversely, if all these determinants vanish, \(U\) is a weak lift.  In the model, \(U_2\) is not a weak lift.
\end{problem}

Here is the statement of the proof that turns these lengths into a Hecke element, in its simplest case: no auxiliary places, no tame elements, one character.  Let \(G\) be finite, \(U_1, \dots, U_n \colon G \to \GL_q(\Oc)\), and \(p \ne 0\) in \(\Oc\).  Let \(\mathcal S \subseteq \Oc^n\) be an \(\Oc\)-subalgebra (componentwise operations) containing the families \((c_k(U_i(g)))_i\) of coefficients of characteristic polynomials, and \(\mathcal J \subseteq \Oc^n\) an \(\mathcal S\)-submodule containing the families \((p \tr U_i(g))_i\).  Let \(\mathcal X \subseteq \{1, \dots, n\}\) contain \(1\), and suppose that, in one basis for all \(i \in \mathcal X\), \(U_i = \rho_i \oplus 1\), and some \(g_i\) has \(\rho_i(g_i) = y_i \cdot 1\) with \(y_i - 1 \in \Oc^\times\); suppose the values of \(\rho_1\) span \(M_{q-1}(\Oc)\).  Let \(b_{\mathrm{main}} \in \mathbb N\) bound the length of the Ribet cocycles for \(\rho_1\) with values in \((\mathrm{Frac}(\Oc)/\Oc)^{q-1}\) normalized at \(g_1\), and \(b_0\) the length of the Ribet cocycles for the trivial character with values in \(\Oc/p\).  Then the families \(\delta_\omega := (\det \sum_g \omega(g)(U_i(g) - 1))_i\) lie in \(\mathcal S\) and vanish on \(\mathcal X\), and some \(\mathfrak f \in \mathcal S\) has
\[
  \mathfrak f\,\delta_\omega \in \mathcal J \text{ for every } \omega, \qquad \ord \mathfrak f_1 \le b_{\mathrm{main}} + b_0 .
\]
\emph{Status}: read from the proof \cite{GL}; it is the specialization of its first-order lemma to a finite group, with no auxiliary places and no tame elements, one trivial character, the whole group as detecting set and separating scalar \(1\).  We take it as given.

\begin{problem}\label{pr:model-element}
In the model take \(\mathcal S = \Tb\), \(p = 3^\nu\) with \(\nu \ge 1\), \(\mathcal J = 3^\nu\Tb\), \(\mathcal X = \{1\}\), \(\rho_1 = \mathrm{sgn}\), \(g_1 = \tau\).  The hypotheses hold with \(b_{\mathrm{main}} = 1\), \(b_0 = 0\); every \(\delta_\omega\) is \((0, d)\) with \(d \in 3\Z_3\), and \(d = 3\) occurs.  The element \(\mathfrak f = (3, 0)\) qualifies, and every \(\mathfrak f \in \Tb\) with \(\mathfrak f\delta_\omega \in \mathcal J\) for all \(\omega\) has \(\ord \mathfrak f_1 \ge 1\).
\end{problem}

So the bound is attained by \(\mathfrak f = (3, 0)\), which vanishes at \(U_2\) and has value \(3\) at the lift; this value generates the congruence ideal \(3\Z_3\) of Problem~\ref{pr:model-congruence}.

\section{The costs and the bound for one operator}\label{sec:costs}

We now use the data of Chapter~\ref{da-chap}, the integer \(\mathcal N\) of Section~\ref{ob-sec:chosen}, and the setting and definitions of Section~\ref{dp-sec:statement}: \(\ell\), \(\iota\), \(\val\), \(\val_+\), coefficient rings \(\Oc\), \(G_E\), \(\cyc\), and \(n_\ell = \val_\ell(\mathcal N)\) and \(\varepsilon_\ell\) of Definition~\ref{dp-def:clearing-bound}; in this sheet \(o_\ell := \val_\ell(X_{\mathrm{odd}})\), so that the first term of \(\varepsilon_\ell\) is \(t\,o_\ell\).  For an isotypic system containing a nonzero lift generator, the trace identity of Definition~\ref{dp-def:main-summand} gives the trace of the normalized representation as the trace of the main summand twisted by \(\hat\vartheta\), plus \(\sum_{0 \le e < t'}\cyc^{-e}\).  Separating these \(t'\) characters costs a power of \(\ell\).

\begin{definition}\label{def:sep}
Put \(h_\ell := 0\) if \(\ell > \half^2\) and \(h_\ell := n_\ell\) otherwise, and \(\sep := \ell^{\half h_\ell}\).  (The proof of the generator lemma supplies \(h_\ell\) as a witness; with this witness the separation scalar is \(\sep\).)
\end{definition}

\begin{problem}\label{pr:generator-sep}
Use: \(p_0 \ne p_1\); \(t = 5r\) with \(r \ge 2\) and \(t \le \half^2\); \(E \subseteq \Q(\mu_{q_1d})\) with \(q_1d \le \half^2\); \(\mu_4 \subseteq E\); each \(p_i\) has a chosen place \(w\) with \(k_w = \F_{p_i}\); and \(p_i^a - 1\) divides \(\mathcal N\) for \(1 \le a \le rt\).  Then some \(g_{\mathrm{cyc}} \in G_E\) has \(\cyc(g_{\mathrm{cyc}})\) generating \(\cyc(G_E)\) topologically and \(\cyc(g_{\mathrm{cyc}})^a - 1 \mid \ell^{h_\ell}\) in \(\Z_\ell\) for \(0 < a < t\).
\end{problem}

\begin{problem}\label{pr:separate}
For the values \(1, 4, 16\) at \(\ell = 3\), the products of differences are \(45, -36, 180\), each of \(3\)-adic order \(2\), and \(9\) times each Lagrange interpolation polynomial has \(3\)-integral coefficients.  In general, with \(g_{\mathrm{cyc}}\) as in Problem~\ref{pr:generator-sep}, \(u := \cyc(g_{\mathrm{cyc}})\), \(r \ge 2\) and \(t' \le t\), every product \(\prod_{i < t',\, i \ne e}(u^e - u^i)\), \(e < t'\), divides \(\sep\).
\end{problem}

Tame ramification costs too.  A relation \((U_i(\delta) - 1)\,o(U_i(\delta)) = 0\) constrains the value at \(\delta\) of a Ribet cocycle for a character trivial at \(\delta\).

\begin{problem}\label{pr:tame}
\begin{enumerate}
\item[(a)] If \((U_i(\delta) - 1)\,o(U_i(\delta)) = 0\) for all \(i\), with \(o \in \Oc[X]\), and \(\xi(\delta) = 1\) for a character \(\xi\), then every Ribet cocycle \(z\) for \(\xi\) has \(o(1)z(\delta) = 0\).  If \((U - 1)(U - \lambda) = 0\), take \(o = \lambda - X\); if \(U^{\mathcal N} = 1\), take \(o = \sum_{i < \mathcal N} X^i\), so \(o(1) = \mathcal N\).
\item[(b)] In the model let \(\delta\) be a \(3\)-cycle.  Then \((X - 1)(1 + X + X^2)\) annihilates \(U_1(\delta)\) and \(U_2(\delta)\), and \(3z(\delta) = 0\) for every Ribet cocycle \(z\) for \(\mathrm{sgn}\).
\item[(c)] Let \(\vartheta\) be auxiliary with ramification at most \(\mathrm{ram}\), choose tame generators \(s_w\) witnessing the bound of Definition~\ref{dp-def:aux-twist}.  At each \(w \in Q_\vartheta\) assign \(o = \hat\vartheta(\delta) - X\) to each of \(\delta = s_w\) and \(\delta = (s_w^c)^{-1}\); at each of the at most \(4f\) places above \(M\) not above \(\ell\) assign \(o = \sum_{0 \le i < \mathcal N}X^i\) to its tame generator and to the inverse of its conjugate.  Then the sum over these indexed terms of \(\val_+(o(1))\) is at most \(2\,\mathrm{ram} + 8f n_\ell\).
\end{enumerate}
\end{problem}

In the general first-order lemma of the proof, \(\Oc\) is a discrete valuation ring, \(G\) is a group, \(p \in \Oc\) is nonzero, and \(U_i \colon G \to \GL_q(\Oc)\) is a finite family of representations.  At every lift index \(i \in \mathcal X\) the matrices are, in common coordinates, \(\rho_i \oplus \bigoplus_{e < t'}\xi_e\), with \(1 \le t' \le q\) and characters \(\xi_e\), \(\xi_0 = 1\).  At a fixed \(i_0 \in \mathcal X\) the values of \(\rho_{i_0}\) span \(M_{q-t'}(\Oc)\); for each \(i \in \mathcal X\) some \(g_i\) has \(\rho_i(g_i) = y_i 1\), \(y_i - 1 \in \Oc^\times\), and \(\xi_e(g_i) = 1\) for all \(e\).  An element \(g_{\mathrm{sep}} \in G\) and a scalar \(\mathsf C \in \Oc\) satisfy \(\prod_{k < t',\, k \ne e}(\xi_e(g_{\mathrm{sep}}) - \xi_k(g_{\mathrm{sep}})) \mid \mathsf C\) for every \(e\).  Let \(J\) be a finite linearly ordered set of \(n_{\mathrm{aux}}\) indices with a nondecreasing assignment \(j \mapsto e_j \in \{0, \dots, t'-1\}\).  For a commutative \(\Oc\)-algebra \(L\), let \(\mathcal S\) be an \(\Oc\)-subalgebra of \(\prod_i L[X_j : j \in J]\) containing the families of variables and of characteristic coefficients, and let \(\mathcal J\) be an \(\mathcal S\)-submodule of this product containing \((p\tr U_i(g))_i\) for every \(g\); coefficients from \(\Oc\) are read as constant polynomials via \(\Oc \to L\).  For each \(j\) there are \(\epsilon_j \in \Oc[X]\) and two pairs \((s_{jb}, \mathrm{Fr}_{jb})\) in \(G\), \(b = 0, 1\), such that for every \(g \in G\) the family of traces of \(\bigl(U_i(s_{jb}) - 1 - X_j(U_i(\mathrm{Fr}_{jb}) - 1)\epsilon_j(U_i(\mathrm{Fr}_{jb}))\bigr)U_i(g)\) lies in \(\mathcal J\), and \(\epsilon_j(\xi_e(\mathrm{Fr}_{jb})) \equiv 1 \pmod p\) for all \(e\).  A finite family \((\delta_w, o_w)\) in \(G \times \Oc[X]\) has \(\xi_e(\delta_w) = 1\) and \((X - 1)o_w(X)\) annihilating every \(U_i(\delta_w)\).  Let \(b_{\mathrm{main}} \in \mathbb N\) and \(b_e \in \mathbb N \cup \{\infty\}\) for \(0 \le e < t'\).  Suppose \(b_{\mathrm{main}}\) bounds the length of the Ribet cocycles for \(\rho_{i_0}\) with values in \((\mathrm{Frac}(\Oc)/\Oc)^{q-t'}\) that vanish at \(g_{i_0}\) and at all \(s_{jb}\).  For each \(e\), every submodule \(\mathcal Z\) of Ribet cocycles for \(\xi_e\) with values in \(\Oc/p\) whose members vanish at all \(\delta_w\) and satisfy \(z(g_{i_0}) = 0 \Rightarrow (\xi_e(g) - 1)z(g_{\mathrm{sep}}) = 0\) for all \(g\), satisfies the triangle cost condition of the proof along the indices of type \(e\), with scalar \(\mathsf C\) and bound \(b_e\), and is detected by one fixed set of \(n_{\mathrm{det}}\) elements (a member vanishing there is zero).  Then every \(\delta_\omega\), \(\omega\) finitely supported, lies in \(\mathcal S\) and vanishes on \(\mathcal X\); one \(\mathfrak f \in \mathcal S\) has \(\mathfrak f\delta_\omega \in \mathcal J\) for all \(\omega\); at every \(i \in \mathcal X\), \(\mathfrak f_i\) has no monomial properly dividing \(\prod_j X_j\); and the coefficient of \(\prod_j X_j\) in \(\mathfrak f_{i_0}\) is the image of some \(c \in \Oc\) with \(\ord c\) at most
\[
  b_{\mathrm{main}} + \sum_{0 \le e < t'} b_e + t'\Bigl(\sum_w \ord o_w(1) + 2n_{\mathrm{det}} \ord \mathsf C\Bigr) + (n_{\mathrm{aux}} + 1)\ord \mathsf C .
\]
For the exponent \(h\) and the element \(g_{\mathrm{cyc}}\) supplied by the generator lemma, the application takes \(\xi_e = \cyc^{-e}\), \(g_{\mathrm{sep}} = g_{\mathrm{cyc}}^{-1}\) and \(\mathsf C = \ell^{\half h}\); at each auxiliary place the two ties pair its tame generator with its Frobenius element, and the inverses of their conjugates with one another.  We write the estimates with the witness \(h = h_\ell\) of Definition~\ref{def:sep}.  In this application \(p = \ell^\nu\).  Over a sufficiently large coefficient ring choose the twisted lattice \(L_{\mathrm{tw}}\) at \(i_0\) with the prescribed Selmer bound, and write \(g_* = g_{i_0}\).  The proof takes \(b_{\mathrm{main}}\) to be the sum of the local dual defects of the Tate duals of \(L_{\mathrm{tw}}\) above \(\ell\) and the length of its crystalline Selmer module with divisible coefficients normalized at \(g_*\), whose normalized length is at most \(\mathrm{sel}\).  For \(0 \le e < t'\), \(b_e\) is the crystalline colength of the Ribet cocycles for \(\cyc^{-e}\) modulo \(\ell^\nu\), plus the length of the module of Definition~\ref{tt-def:shaodd} at \(e\) and \(\nu\) (at \(e = 0\), all everywhere unramified continuous homomorphisms to \(\Oc/\ell^\nu\)), whose normalized length is at most \(o_\ell\).  The normalized sum of the main dual defects and of the character colengths is at most \(r'\loc_\ell\).  Here \(n_{\mathrm{test}}\) is the number of test points plus one for the separating element, \(n_{\mathrm{det}} \le 2(n_{\mathrm{aux}} + n_{\mathrm{test}})\) and \((n_{\mathrm{aux}} + n_{\mathrm{test}})\half h_\ell \le \half^{22}h_\ell\).  No length of a module of even classes enters.

\begin{problem}\label{pr:arith}
Let \(t', h, n_{\mathrm{aux}}, n_{\mathrm{test}}, n_{\mathrm{det}}\) be natural numbers and \(\mathrm{tame}, \mathrm{ram} \in [0, \infty]\), with \(t' \le t\), \(h \le n_\ell\), \((n_{\mathrm{aux}} + n_{\mathrm{test}})\half h \le \half^{22}h\), \(n_{\mathrm{det}} \le 2(n_{\mathrm{aux}} + n_{\mathrm{test}})\) and \(\mathrm{tame} \le 2\,\mathrm{ram} + 8fn_\ell\).  Then
\[
  t'o_\ell + t'\bigl(\mathrm{tame} + 2n_{\mathrm{det}}\half h\bigr) + (n_{\mathrm{aux}} + 1)\half h + 2n_\ell \le 2t'\,\mathrm{ram} + \varepsilon_\ell .
\]
\end{problem}

At auxiliary level the proof realizes the first-order element as a polynomial over the Hecke algebra.  Let \(A\) be the indicator of the systems whose normalized characteristic coefficients are congruent modulo the maximal ideal to those of the distinguished lift at every \(g \in G_E\); multiplying the coefficient of \(\prod_j X_j\) by \(A^{n_{\mathrm{aux}} + 1}\) gives diagonal scalars.  For a nonzero lift generator of type \(c\), the proof projects to type \(c\), passes to smooth forms, applies the diagonal operator and scales by \(\mathcal N^2\).  One factor \(\mathcal N\) makes the type projection of an integral input integral, and the other comes from the norm operator at auxiliary level; their valuation is \(2n_\ell\).  It does this for a basis of lift generators and takes a combination with coefficients integral at \(\iota\).  The result is the theorem of this sheet.

\begin{theorem}\label{thm:eigenvalues}
Let \(\vartheta\) be a twist that is auxiliary with ramification at most \(\mathrm{ram} \in [0, \infty]\), let \(\mathrm{sel} \in [0, \infty)\), let \(\DP\) be a descent pair of ranks \(r' < n'\) and \(t' = n' - r'\), \(e_1\) an element of the group algebra of \(\prod_w \GL_{r'}(k_w)\) over its chosen places, and \(\mathcal F \subseteq \mathcal A(\DP)\) a subspace contained in the span of the lift generators for \((\vartheta, \DP, e_1, \mathcal F)\) with Selmer bound \(\mathrm{sel}\).  Then some \(C_\ell \in \operatorname{End}_{\C} \mathcal F\) has the following properties.
\begin{enumerate}
\item[(a)] For every split place \(w\) such that neither \(w\) nor \(\bar w\) is a place of \(\DP\), all \(T, a \in \C\) and \(y \in \mathcal F\): if the Hecke polynomial operator of \(\DP\) at \(w\) at the value \(T\) sends \(y\) to \(ay\), it sends \(C_\ell y\) to \(aC_\ell y\).
\item[(b)] \(\val_+(\iota\mu) \le \mathrm{sel} + 2t'\,\mathrm{ram} + r'\loc_\ell + \varepsilon_\ell\) for every eigenvalue \(\mu\) of \(C_\ell\).
\item[(c)] \(C_\ell(\mathrm{pr}(y))\) is integral at \(\iota\) for every form \(y\) integral at \(\iota\) and orthogonal to the false lifts for \((\DP, e_1, \mathcal F^\perp)\), where \(\mathrm{pr}\) is the orthogonal projection onto \(\mathcal F\).
\end{enumerate}
\end{theorem}

\emph{Status}: read from the proof \cite{GL}; not proved here.  Problems~\ref{pr:model-element}, \ref{pr:separate}, \ref{pr:tame} and~\ref{pr:arith} model its bound and prove its arithmetic.

\begin{problem}\label{pr:deduce}
Let the hypotheses of Theorem~\ref{dp-thm:denominators} hold.  Use three facts of the proof: if some lift generator is nonzero, then \(1 \le r' \le r\), \(t' \le t\) and \(n' \le 2n_m\) (Definition~\ref{dp-def:lift-setting}); for forms of rank \(n' \le 2n_m\) at a level inside the standard one, the orders of the stabilizers in the class sum divide \(\mathcal N\); and for \(y_1\) integral at \(\iota\) and \(y'\) integral at \(z \mapsto \iota(\bar z)\), each term of the class sum \(\langle y_1, y' \rangle\), without its weight, is integral at \(\iota\).  Deduce that the projection onto \(\mathcal F\) has denominator at most \(\cl_\ell\log\ell\) at \(\iota\) in the sense of Definition~\ref{dp-def:denominator}.
\end{problem}

What remains is the first-order element itself.  The construction ties the rows of the auxiliary places and uses an adjugate identity to turn the relations into a multiplier.  Amitsur's formula puts the coefficients of its scalar polynomial in \(\mathcal S\), which contains the characteristic coefficients and the auxiliary variables \cite{Am} (see \cite[formula (1.5)]{Ch}).  The passage from equal traces to block lattices uses the Brauer--Nesbitt theorem for trace-nondegenerate representations \cite[(30.16)]{CR}.  These steps and the choice of the auxiliary places are taken from the proof here.

\paragraph{What this sheet establishes.}  On \(S_3\) at \(3\) the congruence exponent equals the length of the normalized Ribet cocycles, and the specialization of the first-order bound to this model is attained.  In general the sheet proves the Fitting-ideal step in absolute form, the separation and tame costs, the arithmetic that turns the count into the bound of Theorem~\ref{thm:eigenvalues}, and the deduction of Theorem~\ref{dp-thm:denominators}.  We take Theorem~\ref{thm:eigenvalues}, the first-order lemma and its inputs \(b_{\mathrm{main}}\), \(b_e\) as given, with Chapter~\ref{n3-chap} for the local loss and Theorem~\ref{tt-thm:oddsha} for the odd modules.  Not treated here: the choice of the auxiliary places, and why no length of even classes enters the costs \(b_e\).

\paragraph{In the Lean formalization.}
{\raggedright
In the Lean proof \cite{GL} (\texttt{grh/src/}; \texttt{FO} is \texttt{New/Congruences/FirstOrder}; names without a namespace are in \texttt{New.Selmer}), Theorem~\ref{thm:eigenvalues} is \texttt{exists\_operator\_ellVal\_le\_liftValuationBound} (\texttt{FO/DenominatorBound.lean}), with (a) the predicate \texttt{PreservesHeckeEigenspaces}.  The conclusion of Problem~\ref{pr:deduce} is \texttt{denominatorLe\_clearingBound}, through \texttt{exists\_operator\_det\_clearingBound}; its three facts are \texttt{IsLiftGenerator.exists\_type} with the bounds of \texttt{IsLiftSetting}, \texttt{New.Data.card\_classAut\_dvd\_chosenGroupOrder} and \texttt{New.gramForm\_mem\_embValRing}.  Definition~\ref{def:ribet} is \texttt{ribetCocyclesIn} (\texttt{New/Congruences/LocalLoss/Basic.lean}), and the normalized ones are \texttt{ribetCocyclesVanishingOn} with empty vanishing set; Definition~\ref{def:weak} is the condition of \texttt{IsWeakLift}.  The exponent of Definition~\ref{def:sep} is the choice made in the proof of \texttt{exists\_cyclotomicChar\_pow\_sub\_one\_dvd}, which gives Problem~\ref{pr:generator-sep}; the general product bound of Problem~\ref{pr:separate} follows from \texttt{mul\_prod\_pow\_sub\_pow\_dvd} with multiplier \(1\), \(t'\) characters and \(y = \Delta/2\).  The relative form of Problem~\ref{pr:fitting-step} is \texttt{fittingIdeal.exists\_addVal\_le\_relLength} (\texttt{Lib/General/Algebra/FittingIdealsAndLengths.lean}); the block statement is \texttt{exists\_block\_rows}, applied through \texttt{exists\_tiedRows\_addVal\_det\_le}.  The first-order lemma is \texttt{exists\_isCommonRibetElement\_addVal\_le} (\texttt{FO/RibetLemma.lean}); the proof assembles its inputs in \texttt{ribet\_columns}, with \(b_{\mathrm{main}}\) from \texttt{exists\_length\_le\_loss\_add\_crystallineSelmerLength} and the costs \(b_e\) from \texttt{hasTriangleCost\_of\_isCharacterCocycle}, and obtains the diagonal scalars and their integrality estimate in \texttt{exists\_ellVal\_le\_and\_diagOp\_mem\_ellIntegralForms}; the map for one nonzero generator is constructed in \texttt{exists\_linearMap\_eq\_smul\_mem\_integralForms}.  The cost in Problem~\ref{pr:tame}(c) is the valuation estimate in the proof of \texttt{exists\_tame\_relations}, which also constructs the tame family and proves the annihilation under its hypotheses.  Problem~\ref{pr:arith} is \texttt{mul\_padicValNat\_add\_le\_add\_denominatorError} (\texttt{FO/TraceRelations.lean}).  The construction of the element is \texttt{exists\_traceIdeal\_multiplier\_eq\_sum\_det}, with \texttt{exists\_adjugate\_polynomial} for the adjugate.  The predicates \texttt{IsTraceIdeal}, \texttt{IsCharacterCocycle}, \texttt{HasTriangleCost} and \texttt{IsDetectedBy} occur in the hypotheses of the general lemma; its conclusion gives \(c \in \Oc\) satisfying the displayed valuation bound such that \texttt{IsCommonRibetElement} holds with distinguished coefficient the image of \(c\) in \(L\).  The diagonal scalars are \texttt{New.Patching.BaseLevel.exists\_diagOp\_mem\_ellIntegralForms}, with \texttt{New.Patching.BaseLevel.exists\_indicator\_isCongruent}.  The lift trace is \texttt{HasLiftTrace}; the passage to a lattice in block form uses \texttt{exists\_blockDiagonal\_lattice} and \texttt{New.IsTraceNondeg.exists\_conj\_of\_trace\_eq}; the auxiliary places are chosen by \texttt{exists\_taylorWilesPlaces\_hasRootChain}.
\par}

\section*{Hints}

\subsection*{\S\ref{sec:model}}
\begin{description}
\item[\ref{pr:model-traces}] Compute trace and determinant for one element of each conjugacy class, and compare the coefficient pairs with the pairs congruent modulo \(3\).
\item[\ref{pr:model-congruence}] For the reduction, \(U_2\) modulo \(3\) fixes a nonzero vector.
\end{description}

\subsection*{\S\ref{sec:ribet}}
\begin{description}
\item[\ref{pr:normalize}] Evaluate the coboundary at \(g_*\).  For existence, check that the coboundary of \((y - 1)^{-1}z(g_*)\) obeys the relations: multiply a relation on the right by \(U(g_*)\).
\item[\ref{pr:model-lattice}] Hermite normal form over \(\Z\) of the \(8\)-component vectors.
\item[\ref{pr:model-cocycles}] Problem~\ref{pc-pr:graph-map}: a Ribet cocycle is a linear map on \(\Lc\); the cocycle rule says it kills \((U(g) - \mathrm{sgn}(g))(U(h) - 1)\).  For \(U_1\) alone, \(U_1\) is trivial on the \(3\)-cycles.
\end{description}

\subsection*{\S\ref{sec:fitting}}
\begin{description}
\item[\ref{pr:fitt-length}] Elementary divisors.
\item[\ref{pr:forms}] \(\Phi^{-1}\mathsf N\) contains \(p\Oc^n\); apply elementary divisors to it.
\item[\ref{pr:fitting-step}] (b) Expand a determinant with rows in \(\mathsf N + \mathsf R\mathbf w\) multilinearly; for the forms, use Problem~\ref{pr:forms}.  (c) Noetherian induction on \(\mathsf N\).
\item[\ref{pr:det-kills}] The adjugate.
\item[\ref{pr:model-element}] \(\mathrm{Hom}(S_3, \Z/3^\nu) = 0\).  If \(\mathfrak f\,(0, 3) \in 3^\nu\Tb\) then \(3^\nu \mid \mathfrak f_2\).
\end{description}

\subsection*{\S\ref{sec:costs}}
\begin{description}
\item[\ref{pr:generator-sep}] \(\cyc(G_E)\) is procyclic; take \(g_{\mathrm{cyc}}\) with \(\cyc(g_{\mathrm{cyc}})\) generating it, and compare \(\cyc(g_{\mathrm{cyc}})^a - 1\) with \(\cyc(g)^a - 1\) for one convenient \(g\), using closed subgroups.  If \(\ell > \half^2\), then \(\ell\) is odd and \(E\) is unramified at \(\ell\): use an inertia element at \(\ell\) whose cyclotomic character reduces to a generator of \(\F_\ell^\times\), and Problem~\ref{cp-pr:generator}.  Otherwise use a Frobenius element at a chosen place above a \(p_i \ne \ell\), and Problem~\ref{cp-pr:gl-order}.
\item[\ref{pr:separate}] \(u^e - u^i = u^i(u^{e-i} - 1)\), and \(t - 1 \le \half\).
\item[\ref{pr:tame}] (a) Write \(o(X)(X - 1) = \sum_k a_kX^k(X - 1)\) and use \(z(\delta^k) = kz(\delta)\).  (c) Compare the character at \(s_w\) and at \((s_w^c)^{-1}\), using \(\hat\vartheta(g^c) = \hat\vartheta(g)^{-1}\).
\item[\ref{pr:arith}] Expand and compare term by term, with \(t' \le t\).
\item[\ref{pr:deduce}] Apply Theorem~\ref{thm:eigenvalues} with \(\mathrm{sel} = \mathrm{ram} = \mathrm{tw}_\ell\), and try \(\mathcal N C_\ell\), or the identity if \(\mathcal F = 0\).  Problems~\ref{dp-pr:budget} and~\ref{dp-pr:det-bound}; Problem~\ref{cp-pr:pairing} is the model for the pairing.
\end{description}

{\small
\begin{thebibliography}{GL}
\bibitem[Am]{Am} S.~A.~Amitsur, \emph{On the characteristic polynomial of a sum of matrices}, Linear and Multilinear Algebra \textbf{8} (1980), 177--182.
\bibitem[Ch]{Ch} G.~Chenevier, \emph{The \(p\)-adic analytic space of pseudocharacters of a profinite group and pseudorepresentations over arbitrary rings}, in: Automorphic Forms and Galois Representations, Vol.~1, LMS Lecture Note Ser.\ 414, Cambridge Univ.\ Press, 2014, 221--285.
\bibitem[CR]{CR} C.~W.~Curtis and I.~Reiner, \emph{Representation Theory of Finite Groups and Associative Algebras}, Interscience, 1962.
\bibitem[Ei]{Ei} D.~Eisenbud, \emph{Commutative Algebra with a View Toward Algebraic Geometry}, Grad.\ Texts in Math.\ 150, Springer, 1995.
\bibitem[GL]{GL} L.~Sernau, \emph{GRH for Dirichlet \(L\)-functions, in Lean~4}, a shorter rewriting of the Lean proof, text by Claude, \url{https://github.com/lsernau/grh-lean}, commit \texttt{72fead1}, readable edition \texttt{grh/}.
\end{thebibliography}
}

\end{document}
